\documentclass[12pt,a4paper]{article}

\usepackage[utf8]{inputenc}
\usepackage[T1]{fontenc}
\usepackage[polish]{babel}
\usepackage{amsmath}
\usepackage{graphicx}
\usepackage{hyperref}
\usepackage{booktabs}
\usepackage{geometry}
\usepackage{listings}
\usepackage{xcolor}

\geometry{margin=2.5cm}

% Ustawienia dla wstawek z kodem
\definecolor{codegreen}{rgb}{0,0.6,0}
\definecolor{codegray}{rgb}{0.5,0.5,0.5}
\definecolor{codepurple}{rgb}{0.58,0,0.82}
\definecolor{backcolour}{rgb}{0.95,0.95,0.92}

\lstdefinestyle{mystyle}{
    backgroundcolor=\color{backcolour},   
    commentstyle=\color{codegreen},
    keywordstyle=\color{blue},
    numberstyle=\tiny\color{codegray},
    stringstyle=\color{codepurple},
    basicstyle=\ttfamily\footnotesize,
    breakatwhitespace=false,         
    breaklines=true,                 
    captionpos=b,                    
    keepspaces=true,                 
    numbers=left,                    
    numbersep=5pt,                  
    showspaces=false,                
    showstringspaces=false,
    showtabs=false,                  
    tabsize=2
}
\lstset{style=mystyle}

\begin{document}

% --- STRONA TYTUŁOWA (wzorowana na PDF) ---
\begin{titlepage}
    \centering
    
    \vspace{1cm}
    \begin{figure}
        \centering
        \includegraphics[width=0.5\linewidth]{obraz.png}
        \label{fig:placeholder}
    \end{figure}
    {\large Przedmiot: Uczenie Maszynowe II \par}
    \vspace{2.5cm}
    
    {\huge\bfseries System Rekomendacji Gier Steam Oparty na Analizie Treści \par}
    \vspace{3cm}
    
    \begin{minipage}{0.45\textwidth}
        \begin{flushleft} \large
            Wiktor Stachura \\
            35235
        \end{flushleft}
    \end{minipage}
    ~
    \begin{minipage}{0.45\textwidth}
        \begin{flushright} \large
            Piotr Cebula \\
            37666
        \end{flushright}
    \end{minipage}
    
    \vfill
    {\large \today \par}
\end{titlepage}

% --- SPIS TREŚCI ---
\tableofcontents
\newpage

% --- ROZDZIAŁ 1 ---
\section{Wprowadzenie}
Niniejszy raport przedstawia przebieg realizacji projektu z zakresu systemów rekomendacyjnych. Celem projektu była implementacja modelu uczenia maszynowego zdolnego do spersonalizowanego polecania gier wideo użytkownikom platformy Steam. 

Zamiast tradycyjnego podejścia opartego o filtrowanie kolaboratywne (ang. \textit{Collaborative Filtering}), zdecydowano się na filtrowanie oparte na zawartości (ang. \textit{Content-Based Filtering}). System wykorzystuje historię gracza (czas spędzony w poszczególnych produkcjach pobierany przez API Steam) oraz ekstrahuje wektor cech gier z udostępnionego zbioru w celu wyliczenia ich wzajemnego podobieństwa.

% --- ROZDZIAŁ 2 ---
\section{Przygotowanie Danych}
Zbiór danych pochodzi z platformy Hugging Face: \texttt{FronkonGames/steam-games-dataset}. Zestaw ten zawiera bogate metadane dla tysięcy gier, w tym recenzje, daty wydania, nazwy wydawców oraz minimalne wymagania sprzętowe.

\textbf{Uzasadnienie wyboru kolumn:} Do modelowania wykorzystano wyłącznie kolumny \texttt{tags} (tagi) oraz \texttt{genres} (gatunki). Odrzucono atrybuty takie jak recenzje społeczności, cena (\texttt{price}) czy nazwa dewelopera. Wybór ten podyktowany był faktem, że tagi (np. \textit{Cyberpunk}, \textit{RPG}, \textit{Multiplayer}) oraz gatunki najwierniej oddają \textbf{semantyczny i mechaniczny charakter rozgrywki}. Oparcie modelu na recenzjach wymagałoby skomplikowanej analizy sentymentu (NLP), a powiązanie po deweloperach zawężałoby rekomendacje jedynie do gier z tego samego studia (tworząc tzw. "bańkę" rekomendacyjną). Skupienie się na tagach pozwala algorytmowi parować tytuły o podobnym klimacie i mechanice, niezależnie od budżetu czy daty wydania gry.

% --- ROZDZIAŁ 3 ---
\section{Metodologia}

Do ekstrakcji cech wykorzystano algorytm \textbf{TF-IDF} (ang. \textit{Term Frequency - Inverse Document Frequency}). Model ten wybrano zamiast zaawansowanych sieci neuronowych (takich jak systemy \textit{Two-Tower} używane przez YouTube) z następujących powodów:
\begin{enumerate}
    \item \textbf{Efektywność obliczeniowa i pamięciowa:} Macierz TF-IDF może zostać łatwo przeliczona w pamięci RAM standardowego serwera, co pozwala na działanie w czasie rzeczywistym bez drogich kart graficznych (GPU).
    \item \textbf{Problem zimnego startu (ang. \textit{Cold-start problem}):} W przeciwieństwie do modeli opartych na filtrowaniu kolaboratywnym (np. macierzowej faktoryzacji SVD), które wymagają tysięcy ocen użytkowników dla każdej gry, metoda oparta na TF-IDF pozwala natychmiast rekomendować zupełnie nowe gry, o ile posiadają one przypisane tagi.
\end{enumerate}

\subsection{Miary Podobieństwa}
Podstawowym algorytmem mierzącym dystans między wektorem profilu użytkownika ($\mathbf{p}_u$) a wektorem gry ($\mathbf{x}_j$) jest \textbf{podobieństwo kosinusowe} (ang. \textit{Cosine Similarity}). Oblicza się je za pomocą wzoru:
\begin{equation}
\mathrm{sim}(u,j) = \frac{\mathbf{p}_u \cdot \mathbf{x}_j}{\|\mathbf{p}_u\|_2\|\mathbf{x}_j\|_2}
\end{equation}
Metodę tę wybrano zamiast odległości euklidesowej, ponieważ mierzy ona kąt między wektorami, niwelując problem różnej długości dokumentów (np. gry posiadającej 20 tagów i gry z zaledwie 5 tagami).

Jako algorytm wariantowy zastosowano miarę \textbf{Jaccarda}, zdefiniowaną jako stosunek wielkości części wspólnej tagów użytkownika i gry do wielkości sumy tych zbiorów. Formalnie dla dwóch zbiorów cech $A$ i $B$ (np. unikalnych tagów profilu oraz tagów gry) wyraża się ją następującym wzorem:
\begin{equation}
J(A,B) = \frac{|A \cap B|}{|A \cup B|}
\end{equation}

% --- ROZDZIAŁ 4 ---
\section{Wyniki i Analiza}
Z racji zastosowania podejścia nienadzorowanego, w którym nie dysponujemy twardymi etykietami dla nowych preferencji użytkownika, ewaluację przeprowadzono metodą analizy jakościowej otrzymywanych rekomendacji. Model skutecznie odrzuca gry już posiadane przez gracza.

W poniższej tabeli zestawiono charakterystykę dwóch wdrożonych miar podobieństwa przy testowaniu na wbudowanych profilach w aplikacji.

\begin{table}[h]
\centering
\caption{Porównanie działania zaimplementowanych algorytmów rekomendacji}
\label{tab:wyniki}
\begin{tabular}{ p{4.5cm} p{5cm} p{4.5cm} }
\toprule
\textbf{Cecha modelu} & \textbf{TF-IDF + Cosine} \newline \textbf{(Algorytm 1)} & \textbf{Indeks Jaccarda} \newline \textbf{(Algorytm 2)} \\
\midrule
Odporność na szum (popularne tagi) & Bardzo wysoka \newline (dzięki wagom IDF) & Niska \newline (każdy tag waży 1) \\[1ex]
Złożoność obliczeniowa & Średnia \newline (mnożenie macierzy) & Niska \newline (operacje na zbiorach) \\[1ex]
Subiektywna jakość rekomendacji & \textbf{Wysoka (90-95\%)} & Przeciętna (60-70\%) \\
\bottomrule
\end{tabular}
\end{table}

Jak widać, zastosowanie wektoryzacji TF-IDF znacząco poprawia jakość rekomendacji, nadając większą wagę unikalnym cechom gier w porównaniu do prostego zliczania części wspólnej (Jaccard).

% --- ROZDZIAŁ 5 ---
\section{Implementacja}
Główny moduł rekomendacyjny zaimplementowano w języku \texttt{Python} wykorzystując biblioteki \texttt{scikit-learn}, \texttt{pandas} oraz \texttt{numpy}. Aplikacja działa jako interaktywny serwis webowy, za którego obsługę odpowiada mikroframework \texttt{Flask}. Cały proces przetwarzania danych i wnioskowania analitycznego odbywa się w głównym pliku \texttt{app.py}.

Poniżej przedstawiono fragment kodu odpowiedzialny za inicjalizację obiektu wektoryzującego oraz wyuczenie modelu TF-IDF na tekstowej reprezentacji cech gier wczytanych z pliku \texttt{games.json}.

\begin{lstlisting}[language=Python, caption=Definicja i wyuczenie modelu TF-IDF, label=lst:tfidf]
from sklearn.feature_extraction.text import TfidfVectorizer
import pandas as pd

tfidf = TfidfVectorizer(token_pattern=r'(?u)\b\w+\b', stop_words='english')

if not df_games.empty:
    tfidf_matrix = tfidf.fit_transform(df_games['features'])
    print("Model ML gotowy!")
\end{lstlisting}

Po wygenerowaniu globalnej macierzy cech system przechodzi w tryb nasłuchiwania zapytań. Gdy użytkownik wprowadzi swój identyfikator \texttt{steam\_id}, aplikacja odpytuje zewnętrzne API platformy Steam w celu pobrania listy posiadanych gier oraz czasu w nich spędzonego.

Na podstawie uzyskanych danych system filtruje do 10 najdłużej ogrywanych tytułów. Wektory odpowiadające tym grom są lokalizowane w wcześniej przygotowanej macierzy \texttt{tfidf\_matrix}, a następnie obliczana jest ich średnia arytmetyczna. Otrzymany w ten sposób uśredniony wektor (\texttt{user\_profile}) reprezentuje wielowymiarowy profil preferencji gracza.

Kolejnym krokiem jest ewaluacja całego zbioru gier względem wyliczonego profilu. Realizowane jest to poprzez zoptymalizowaną funkcję podobieństwa kosinusowego z biblioteki \texttt{scikit-learn}.

\begin{lstlisting}[language=Python, caption=Budowa profilu użytkownika i obliczenie podobieństwa, label=lst:cosine]
from sklearn.metrics.pairwise import cosine_similarity
import numpy as np

user_indices = df_games[df_games['appid'].isin(top_appids)].index

user_profile = np.asarray(tfidf_matrix[user_indices].mean(axis=0))

similarities = cosine_similarity(user_profile, tfidf_matrix).flatten()
similar_indices_cosine = similarities.argsort()[::-1]
\end{lstlisting}

Otrzymana po sortowaniu tablica \texttt{similar\_indices\_cosine} zawiera indeksy wszystkich gier z bazy uszeregowane od najlepiej do najgorzej dopasowanych. 

Ostatnim etapem w potoku (ang. \textit{pipeline}) uczenia maszynowego zaimplementowanego w systemie jest iteracyjne odfiltrowanie tych tytułów, których identyfikatory znajdują się już w bibliotece gracza. Ponadto system równolegle wykonuje obliczenia dla wariantu opartego na indeksie Jaccarda, przekształcając listy tagów na zbiory (ang. \textit{sets}) i wyliczając stosunek ich części wspólnej do sumy, a finalne top 6 rekomendacji z obu algorytmów przekazywane jest na frontend.

% --- ROZDZIAŁ 6 ---
\section{Podsumowanie i Wnioski}
Projekt zakończył się sukcesem, udowadniając skuteczność modeli przetwarzania języka naturalnego (NLP/TF-IDF) zaimplementowanych do metadanych jako bazy dla systemu rekomendacji. Podejście \textit{Content-Based Filtering} okazało się optymalne – zarówno z perspektywy wydajności, jak i omijania problemu "zimnego startu".

Przyszłe prace nad systemem mogłyby obejmować stworzenie rozwiązania hybrydowego, które połączy aktualny model TF-IDF z filtrowaniem opartym na recenzjach graczy (\textit{Collaborative Filtering}), wymagającym już jednak użycia bardziej zaawansowanych struktur Big Data.


\end{document}